% Exported from Math++. Compile with pdflatex.
\documentclass[11pt,a4paper]{article}
\usepackage[T1]{fontenc}
\usepackage[utf8]{inputenc}
\usepackage{lmodern}
\usepackage{amsmath,amssymb,amsthm,mathtools}
\usepackage{graphicx,booktabs,xcolor}
\usepackage[hidelinks]{hyperref}
\usepackage[margin=25mm]{geometry}
\theoremstyle{plain}
\newtheorem{theorem}{Theorem}
\newtheorem{lemma}[theorem]{Lemma}
\newtheorem{corollary}[theorem]{Corollary}
\newtheorem{proposition}[theorem]{Proposition}
\newtheorem{conjecture}[theorem]{Conjecture}
\theoremstyle{definition}
\newtheorem{definition}[theorem]{Definition}
\newtheorem{example}[theorem]{Example}
\theoremstyle{remark}
\newtheorem*{remark}{Remark}
\setlength{\parindent}{0pt}\setlength{\parskip}{0.6em}
\begin{document}
\section*{Using files together}

A document can use what other files define. These two names come from the
\texttt{lib} folder next to this file.

Using $g, A$ from [[lib/physics.mpp]].

\begin{equation*}
  g = 9.81
\end{equation*}
\begin{equation*}
  A(3) = 9 \pi
\end{equation*}
Everything from another file, with \texttt{![[$\ldots$]]}:

\textit{Imports everything from lib/shapes.mpp.}

\begin{equation*}
  P(3) = 6 \pi
\end{equation*}
\end{document}
